\documentclass[10pt,a4paper]{amsart}
\usepackage[margin=19mm]{geometry}
\usepackage{amsmath,amssymb,mathtools}
\usepackage[scr=rsfs,cal=euler]{mathalfa}
\usepackage{xfrac}
\usepackage[unicode,hidelinks,bookmarks=false]{hyperref}
\newcommand{\ie}{i.e.\ }
\newcommand{\cf}{cf.\ }
\newcommand{\resp}{respectively\ }
\newcommand{\Subsection}{\subsection}
\providecommand{\nobreakdash}{\nobreak\hbox{-}}
\setlength{\emergencystretch}{2em}
\multlinegap0pt
\usepackage{amssymb}
\usepackage{mathtools}
\usepackage{dsfont}
\newtheorem{theorem}{Theorem}
\newtheorem{lemma}{Lemma}
\def\E{\mathbb{E}}
\def\N{\mathbb{N}}
\def\R{\mathbb{R}}
\def\cov{\text{Cov}}
\def\var{\text{Var}}
\title{Limit theorems for the distance \\of random points in $l_p^n$-balls}
\author{David Alonso-Guti\'errez, Javier Mart\'in Go\~ni, Joscha Prochno}
\date{Source: arXiv:2512.24367v1 (CC BY 4.0)}
\begin{document}
\maketitle

\noindent\emph{Source and licence.} Limit theorems for the distance \\of random points in $l_p^n$-balls. Version arXiv:2512.24367v1 (CC BY 4.0), distributed by arXiv under the Creative Commons Attribution 4.0 International licence, http://creativecommons.org/licenses/by/4.0/, as recorded in the arXiv licence metadata for this exact version and shown on the abstract page https://arxiv.org/abs/2512.24367v1. Full licence notice: \texttt{licenses/arxiv-2512.24367-CC-BY-4.0.txt}.\par\smallskip

\noindent\textit{Editorial scope.} Counted unit: the article's theorem Thm:CLTforlpballs (source0.tex lines 184--194). The article's complete original proof of it is delivered with it (source0.tex lines 776--883, one \texttt{proof} environment of 108 lines, which the article places away from the statement). No mathematics is written, repaired or re-derived: the statement and the proof are the article's own lines, in the article's own notation, and no retained line is substituted or re-flowed.  Retained uncounted as the source-local context the proof needs: lines 345--352; lines 719--727.  Nothing was omitted from any retained span, so the package carries no omission record.  No retained line contains a citation, so the wrapper carries no bibliography and no bibliography entry.  No retained line contains a figure environment, an image-inclusion command or drawing code, so no image is delivered and the package declares no asset dependency.  Editorial formatting only, in addition: an AMS \texttt{amsart} preamble that restates the article's own declarations and macros used by the retained lines, namely the environments \texttt{theorem}, \texttt{lemma} on separate counters, together with the standard packages those lines need.  The retained environments print in this document as Theorem 1, Lemma 1 in order of appearance, whereas the article prints the same items with the numbering of its own full document.  The complete native source of the article as distributed in the arXiv e-print for this exact version is bundled unchanged as \texttt{sources/191889/source0.tex}.
\begin{lemma}[Delta method]
\label{DeltaMethod}
Let $(X_n)_{n\in\N}$ be a sequence of $k-$dimensional random vectors such that $\sqrt{n}(X_n - \mu)$ converges in distribution to a centered Gaussian random vector with covariance matrix $\Sigma$, and let $g:\mathbb{R}^k \to \R^d$ be continuously differentiable at $\mu$ with Jacobian matrix $J_g$. Then,
\begin{align*}
\sqrt{n} (g(X_n) - g(\mu) ) \overset{d}{\longrightarrow} N\left( 0 , J_g \Sigma J_g^T \right)
\end{align*}
provided $J_g \Sigma J_g^T$ is positive definite.
\end{lemma}

\begin{lemma}
\label{LemEquivalentCLTBoundarypball}
Let $1\leq p <\infty$ and let $g_i, g_i'$ with $i=1,...,n$, be independent $p-$generalized Gaussian random variables. Then,
\begin{align*}
\frac{n^{1/p - 1/2}}{\left(\sum _{i=1}^n \left| g_i  \right|^p \right)^{1/p}} \left\Vert   \left( g_i - 
 \frac{\left(\sum _{i=1}^n \left| g_i  \right|^p  \right)^{1/p} }{\left(\sum _{i=1}^n \left| g_i'  \right|^p \right)^{1/p}} g_i' \right)_{i=1}^n \right\Vert _2 \overset{d}{\sim} \frac{n^{1/p - 1/2}}{\left(\sum _{i=1}^n \left| g_i  \right|^p \right)^{1/p}} \left\Vert   \left( g_i - 
   g_i' \right)_{i=1}^n \right\Vert _2 .
\end{align*}
\end{lemma}

\begin{theorem}
\label{Thm:CLTforlpballs}
Let $1\leq p <\infty$, and let $X^{(n)},Y^{(n)} \sim \text{Unif}( \partial B_p^n)$ or $ X^{(n)},Y^{(n)} \sim \text{Unif}(B_p^n)$ be independent random vectors. Then,
\begin{align*}
\sqrt{n} \left( n^{1/p - 1/2} \Vert X^{(n)} - Y^{(n)} \Vert _2 -  \sqrt{2} \sqrt{\frac{\Gamma (3/p)}{\Gamma(1/p)}} p^{1/p}  \right) \overset{d}{\longrightarrow} N(0, \sigma_p ^2),
\end{align*}
where $\sigma_p ^2>0$ is given by
\begin{align*}
\sigma_p ^2 = \frac{p^{2/p} \left( M_p(4) + ((M_p(2))^2 \right)}{4 M_p(2)}.
\end{align*}
\end{theorem}

\begin{proof}[Proof of Theorem \ref{Thm:CLTforlpballs} (boundary case)]
Let us define the function $f: \R^2 \to \R$ given by
\begin{align*}
f(x,y) =  x^{1/2} \frac{1}{y^{1/p}}.
\end{align*} 
First, note that
\begin{align*}
n^{1/p - 1/2} \frac{1}{\Vert G \Vert _p} \Vert G - G' \Vert _2 = f \left( \frac{1}{n} \sum _{i=1}^n \left( \left| g_i - g_i' \right|^2 , |g_i|^p \right)  \right).
\end{align*}
Take $S_{1,n} = \frac{1}{n} \sum _{i=1}^n  \left| g_i - g_i' \right|^2$ and $S_{2,n} = \frac{1}{n} \sum _{i=1}^n  \left| g_i  \right|^p$. Thus, by Lemma \ref{LemEquivalentCLTBoundarypball} we have that 
\begin{align*}
n^{1/p - 1/2} \Vert X^{(n)} - Y^{(n)} \Vert _2 \overset{d}{\sim} f\left( (S_{1,n} , S_{2,n}) \right).
\end{align*}
By the classical CLT, $\left( \sqrt{n} ( S_{1,n} - \E [S_{1,n}]), \sqrt{n} ( S_{2,n} - \E [S_{2,n}]) \right)$ converges in distribution to $N(0, \Sigma )$, where $\Sigma$ is the covariance matrix of $(S_{1,n},S_{2,n})$. Thus, applying the Delta method (Lemma \ref{DeltaMethod}) to $X_n = (S_{1,n} , S_{2,n})$ and the function $f$, we have that 
\begin{align*}
\sqrt{n} \left( f(X_n) - f\left( \E [X_n] \right) \right) \overset{d}{\to} N \left(0, J_f \Sigma J_f^T \right),
\end{align*}
where $J_f$ is the Jacobian matrix of $f$ and $\E[X_n] = (2M_p(2) , M_p(p))$.

Let us compute the covariance matrix 
\[
   \Sigma = 
  \left[ {\begin{array}{cc}
   \var[|g_1 - g_1'|^2] & \cov[|g_1 - g_1'|^2 , |g_1|^p] \\
   \cov[|g_1 - g_1'|^2 , |g_1|^p] &  \var[|g_1|^p] \\
  \end{array} } \right].
\]
For the first term,
\begin{align*}
\var \left[ \left| g_i - g_i' \right|^2   \right] & = \var \left[ | g_1 |^2 + | g_1' |^2
- 2 g_1 g_1'   \right]
 \\
 & = \E \left[ \left( | g_1 |^2 + | g_1' |^2 - 2 g_1 g_1' \right) ^2 \right] -
\left( \E \left[ | g_1 |^2 + | g_1' |^2 - 2 g_1 g_1'   \right] \right)^2
 \\
 & = \E \left[ |g_1|^4 + |g_1'|^4 + 6 |g_1|^2 |g_1'|^2 -4|g_1|^2 g_1g_1' -4|g_1'|^2 g_1g_1'   \right]
 - \left( 2 \E \left[ |g_1|^2  \right] \right)^2
 \\
 & = 2 \E \left[ |g_1|^4 \right] + 6 E\left[ |g_1|^2 \right] \E \left[ |g_1'|^2 \right] - \left( 2 \E \left[ |g_1|^2  \right] \right)^2
 \\
 & = 2 \E \left[ |g_1|^4 \right] + 2 \left( \E \left[ |g_1|^2  \right] \right)^2
 \\
 & = 2 M_p(4) + 2 (M_p(2))^2.
\end{align*}
For the fourth term,
\begin{align*}
\var \left[ \left| g_i\right|^p   \right] = \E \left[  | g_1 |^{2p} \right] -
\left( \E \left[ | g_1 |^p   \right] \right)^2 & = \frac{\Gamma \left( \frac{2p + 1}{p} \right)}{\Gamma \left( \frac{1}{p} \right)} - \left( \frac{\Gamma \left( \frac{p+1}{p} \right)}{\Gamma \left( \frac{1}{p} \right)} \right)^2 
 \\
  & = \frac{1}{p} = M_p(p).
\end{align*}
For the remaining terms, note that
\begin{align*}
\cov[\left| g_i - g_i'\right|^2 , \left| g_i\right|^p ] & = \E [\left| g_i - g_i'\right|^2 \left| g_i\right|^p  ] - \E [\left| g_i - g_i'\right|^2  ] \E [ \left| g_i\right|^p  ].
\end{align*}
Let us compute both summands. For the second one,
\begin{align*}
\E [\left| g_i - g_i'\right|^2] = 2\E [\left| g_1 \right|^2] = 2 \frac{\Gamma \left( \frac{3}{p} \right)}{\Gamma \left( \frac{1}{p} \right)}.
\end{align*}
and
\begin{align*}
\E [\left| g_i \right|^p] = \frac{\Gamma \left( \frac{p+1}{p} \right)}{\Gamma \left( \frac{1}{p} \right)} = \frac{1}{p}.
\end{align*}
For the first one,
\begin{align*}
\E [ \left| g_i - g_i'\right|^2  \left| g_i \right|^p] & = \E[ |g_1|^{p+2} + |g_1|^p|g_1'|^2 - 2g_1 g_1' |g_1|^p ]
 \\
  & =  \E[ |g_1|^{p+2}] +  \E[ |g_1|^{p}]  \E[ |g_1'|^{2}]
 \\
  & = \frac{\Gamma \left( \frac{p+3}{p} \right)}{\Gamma \left( \frac{1}{p} \right)} + \frac{\Gamma \left( \frac{p+1}{p} \right)}{\Gamma \left( \frac{1}{p} \right)} \frac{\Gamma \left( \frac{3}{p} \right)}{\Gamma \left( \frac{1}{p} \right)} = \frac{4\Gamma \left( \frac{p+3}{p} \right)}{3\Gamma \left( \frac{1}{p} \right)} .
\end{align*}
Therefore,
\begin{align*}
\cov[\left| g_i - g_i'\right|^2 , \left| g_i\right|^p ] & = \frac{4\Gamma \left( \frac{p+3}{p} \right)}{3\Gamma \left( \frac{1}{p} \right)} - 2 \frac{\Gamma \left( \frac{3}{p} \right)}{\Gamma \left( \frac{1}{p} \right)} \frac{1}{p}
 \\
 & = \frac{2 \Gamma \left( \frac{p+3}{p} \right)}{3\Gamma \left( \frac{1}{p} \right)} = \frac{2}{3} M_{p} (p+2).
\end{align*}
Thus, we obtain the covariance matrix
\[
   \Sigma = 
  \left[ {\begin{array}{cc}
   2 M_p(4) + 2 (M_p(2))^2 & \frac{2}{3} M_{p} (p+2) \\
   \frac{2}{3} M_{p} (p+2)  &  M_p(p) \\
  \end{array} } \right].
\]
For the last step, let us compute the Jacobian of $f: \R^2 \to \R$ given by
\begin{align*}
f(x,y) =  x^{1/2} \frac{1}{y^{1/p}}  .
\end{align*} 
It is clear that 
\[
   J_f (x,y) = 
  \left[ {\begin{array}{cc}
   \frac{1}{2\sqrt{x}} \frac{1}{y^{1/p}}  & \frac{-\sqrt{2}}{p} \frac{\sqrt{x}}{y^{1+ 1/p}}  \\
  \end{array} } \right] .
\]
Thus, the Jacobian matrix of $f$ evaluated at $\E[X_n] = \left( 2 M_p(2) , 1/p \right)$ is
\[
   J_g = 
  \left[ {\begin{array}{cc}
   \frac{p^{1/p} }{2 \sqrt{2} \sqrt{M_p(2)}}  & - \sqrt{2}  p^{1/p} \sqrt{M_p(2)} \\
  \end{array} } \right] .
\]
Then, applying elementary computations, one can check that the vaue of $\sigma_p ^2$ is given by
\begin{align*}
\sigma_p ^2 = J_g \Sigma J_g ^T = \frac{p^{2/p} \left( M_p(4) + (M_p(2))^2 \right)}{4 M_p(2)}.
\end{align*}
\end{proof}

\end{document}
